\section{Simulation Validation and Sim-to-Real Gap}\label{sec:sim-validation}

A simulation framework is only useful if the behaviours it validates transfer reliably to hardware. This section examines what the multi-fidelity simulation (Section~\ref{sec:simulation}) proved, what it could not test, and the quantitative evidence underpinning confidence in the sim-to-real transfer.

% ─────────────────────────────────────────────────────────────────────
\subsection{Three Simulation Levels and What Each Validated}\label{sec:sim-levels-val}

Each simulation level targets a different subset of the system. Table~\ref{tab:sim-val-matrix} summarises the coverage.

\begin{table}[H]
\centering
\caption{Validation coverage by simulation level. \checkmark = validated, $\sim$ = partially validated, --- = not testable.}
\label{tab:sim-val-matrix}
\small
\begin{tabular}{@{}lcccc@{}}
\toprule
\textbf{System Property} & \textbf{Interactive} & \textbf{SITL + Sim Cam} & \textbf{SITL + Webcam} & \textbf{Video Replay} \\
\midrule
State machine logic (19 states, 32 transitions) & \checkmark & \checkmark & \checkmark & --- \\
MAVLink command interface (arm, takeoff, guided, land) & --- & \checkmark & \checkmark & --- \\
Geofence enforcement and SSSI avoidance & $\sim$ & \checkmark & \checkmark & --- \\
Lawnmower pattern coverage & \checkmark & \checkmark & \checkmark & --- \\
GPS estimation pipeline (pixel $\to$ WGS\,84) & \checkmark & \checkmark & \checkmark & \checkmark \\
Multi-target clustering and classification & \checkmark & --- & --- & $\sim$ \\
Visual servoing convergence & \checkmark & \checkmark & $\sim$ & --- \\
Operator verification workflow & \checkmark & \checkmark & \checkmark & --- \\
Detection under real lighting/texture & --- & --- & $\sim$ & \checkmark \\
Motion blur and vibration effects & --- & --- & --- & $\sim$ \\
Real GPS noise (${\sim}2$--3\,m) & --- & $\sim$ & $\sim$ & \checkmark \\
Wind disturbance on flight stability & --- & --- & --- & --- \\
\bottomrule
\end{tabular}
\end{table}

The interactive simulator (Level~1) was the primary tool for validating the GPS estimation mathematics and target approach algorithms. Over several hundred runs with configurable error injection, the three concurrent estimators---rolling average, cumulative average, and Kalman filter---were benchmarked against ground truth. The SITL level (Levels~2--3) validated the complete MAVLink command pipeline against ArduPilot's C++ firmware, confirming that arming sequences, guided waypoint navigation, mode transitions, and landing commands executed correctly with realistic vehicle dynamics. Video replay (Level~4) bridged the gap to real conditions by exercising the detection model on genuine aerial imagery from a DJI flight.

% ─────────────────────────────────────────────────────────────────────
\subsection{Error Injection as a Validation Tool}\label{sec:error-injection-val}

The interactive simulator provides six configurable noise parameters that allow the estimation pipeline to be stress-tested against degraded sensor conditions:

\begin{enumerate}
    \item \textbf{GPS drift} ($\sigma = 0$--5\,m): Ornstein--Uhlenbeck random walk matching real receiver autocorrelation. At $\sigma = 3$\,m (typical open-sky conditions), the Kalman filter converges to sub-metre accuracy after 15--20 observations, while the rolling average achieves ${\sim}1.5$\,m.

    \item \textbf{Camera shake} (0--15\,px): altitude-dependent angular jitter modelling motor vibration. Detection confidence degrades below 0.5 at shake levels above 10\,px (equivalent to ${\sim}40$\,cm ground displacement at 10\,m altitude), indicating the boundary at which the detector's performance is compromised.

    \item \textbf{Inference throttle} (1--30\,FPS): limiting CV throughput to the Pi's measured 4\,FPS reveals that at a search speed of 8\,m/s and 35\,m altitude, the camera footprint advances ${\sim}1.67$\,m between frames---well within the ${\sim}32$\,m ground footprint width---providing approximately 14 frames of target visibility per flyover pass.

    \item \textbf{Altitude noise}, \textbf{yaw jitter}, and \textbf{FOV calibration error}: compound noise conditions that stress the coordinate transform under worst-case assumptions. A 10\% FOV error produces a proportional 10\% bias in all GPS estimates, confirming the importance of the bench calibration (Section~\ref{sec:fov-results}).
\end{enumerate}

Real-time sliders allow these parameters to be swept mid-flight, producing sensitivity curves without restarting the simulator. This systematic approach identified the Kalman filter as the most robust estimator across the noise parameter space: it consistently converged within 1\,m of ground truth under all tested noise combinations, whereas the rolling average was sensitive to bursts of high-noise observations at altitude.

% ─────────────────────────────────────────────────────────────────────
\subsection{DJI Video Replay as a Sim-to-Real Proxy}\label{sec:dji-proxy}

Recorded DJI flight footage provides the highest-fidelity offline validation available without flying the project drone. The 30\,fps video ($1456 \times 1088$ cropped) was paired frame-by-frame with SRT telemetry (GPS, altitude, heading) and processed through the identical detection and estimation pipeline used in all other simulation levels. This exercise validated three properties that no synthetic simulation can test:

\begin{itemize}
    \item \textbf{Detection under real conditions.} The retrained YOLOv8n model (Section~\ref{sec:cv}) detected the dummy in 87\% of frames where the target occupied more than $20 \times 20$\,px, with a mean confidence of 0.82. Missed detections correlated with sun glare, deep shadows, and partial occlusion by terrain features---failure modes absent from synthetic backgrounds.

    \item \textbf{GPS estimation accuracy.} Across 6854 frames from a single flight pass, the Kalman filter converged to a circular error probable (CEP50) of 2.3\,m, with a maximum single-observation error of 16.5\,m. The CEP50 is consistent with the compound uncertainty of the GPS receiver (${\sim}2$\,m), FOV calibration (${\sim}5$\%), and the discovered 100--200\,ms GPS timing lag that introduces a speed-dependent directional bias of ${\sim}1.6$\,m at 8\,m/s.

    \item \textbf{GPS timing lag discovery.} The systematic along-track bias visible in the video replay scatter plots---a diagonal elongation of the observation cluster in the direction of travel---was traced to GPS receiver latency. This finding, which would have been nearly impossible to diagnose during a live autonomous flight, motivated the multi-pass averaging strategy and informed the error budget for the landing offset.
\end{itemize}

% ─────────────────────────────────────────────────────────────────────
\subsection{Validation Results Summary}\label{sec:val-results}

Table~\ref{tab:sim-val-results} summarises the quantitative validation outcomes across all simulation levels.

\begin{table}[H]
\centering
\caption{Quantitative simulation validation results.}
\label{tab:sim-val-results}
\small
\begin{tabular}{@{}lll@{}}
\toprule
\textbf{Metric} & \textbf{Simulation Result} & \textbf{Real / Video Result} \\
\midrule
State machine end-to-end (19 states) & All transitions correct & Bench-verified (no flight) \\
Search pattern coverage & 100\% of polygon area & Pattern executed in SITL \\
GPS estimate convergence (Kalman, 20+ obs) & $< 0.5$\,m from ground truth & CEP50 = 2.3\,m (DJI video) \\
Detection rate (target $> 20 \times 20$\,px) & 100\% (synthetic background) & 87\% (real video) \\
Mean detection confidence & 0.97 (synthetic) & 0.82 (real video) \\
Inference latency (Pi 5, TFLite) & N/A (laptop) & 207\,ms (4.8\,FPS) \\
Offset landing accuracy & $< 1$\,m (no wind, no GPS noise) & Not yet flight-tested \\
\bottomrule
\end{tabular}
\end{table}

% ─────────────────────────────────────────────────────────────────────
\subsection{Known Sim-to-Real Gaps}\label{sec:known-gaps}

Despite the multi-fidelity approach, several properties remain untestable in simulation and represent the primary sources of residual risk:

\begin{itemize}
    \item \textbf{Wind and turbulence.} No simulation level models wind effects on the airframe. Wind-induced position drift during visual servoing and GPS lock phases will degrade estimation accuracy and may trigger geofence warnings near boundary edges. SITL offers a wind parameter, but the interaction with camera stabilisation and detection reliability is unvalidated.

    \item \textbf{Motion blur at speed.} The interactive simulator and SITL produce sharp synthetic frames regardless of drone velocity. Real flight at 8\,m/s with a 200\,ms exposure could produce motion blur of ${\sim}1.6$\,m ground equivalent, potentially reducing detection confidence. The DJI video provides partial evidence (recorded at a different airspeed and camera system), but the IMX296 global shutter should mitigate this for the project hardware.

    \item \textbf{Real GPS accuracy.} SITL GPS noise (${\sim}0.5$\,m) is optimistic. Real open-sky receivers typically exhibit 2--3\,m CEP, with occasional multipath excursions to 5--10\,m near buildings or terrain. The DJI video analysis (CEP50 = 2.3\,m) provides the best available estimate, but was recorded with a different GPS receiver.

    \item \textbf{RF interference and MAVLink dropouts.} The simulation assumes a perfect MAVLink link. In the field, WiFi, radio control, and telemetry transmitters share the 2.4\,GHz band, potentially causing packet loss or latency spikes. The link-lost timeout (5\,s $\to$ RTL) is implemented but has not been tested under real interference conditions.

    \item \textbf{Sunlight and thermal effects.} Detection performance under direct sunlight, long shadows, and thermal shimmer is only partially captured by the DJI video replay, which was recorded under overcast conditions. The model has not been validated against harsh midday lighting or low sun angles.
\end{itemize}

% ─────────────────────────────────────────────────────────────────────
\subsection{Confidence Assessment}\label{sec:confidence}

Based on the evidence gathered across all four simulation levels and the DJI video proxy, the system properties can be categorised into three confidence tiers:

\textbf{High confidence} (validated at multiple levels with consistent results): the finite state machine and all 32 transitions; the MAVLink command interface (arm, takeoff, guided navigation, land); the lawnmower search pattern geometry; the GPS estimation mathematics and Kalman filter convergence; the operator verification workflow; and geofence enforcement including SSSI avoidance.

\textbf{Moderate confidence} (validated in simulation but with known real-world degradation): detection reliability at altitude (87\% on real video vs 100\% synthetic); GPS estimation accuracy (CEP50 = 2.3\,m on video, sub-metre in simulation); inference throughput on the Pi (4.8\,FPS measured on bench, but not under thermal load during sustained flight).

\textbf{Low confidence} (untested or only indirectly tested): offset landing accuracy under wind; visual servoing stability in gusty conditions; detection performance under harsh lighting; system behaviour under RF interference or GPS multipath; and thermal throttling during extended missions. These properties require validation through the progressive flight testing framework (Section~\ref{sec:five-tier}), advancing from passive observation to full autonomy only after each preceding tier passes.
